
\subsubsection*{3.3.2 Corollary}
\phantomsection
\label{sga3:VIIB:3.3.2}

\noindent\textbf{Corollary 3.3.2.}
{\itshape If \(k\) is a field of characteristic \(0\), the
functor
\[
  \mathbf L\longmapsto\mathscr G(\mathbf L)
\]
is an equivalence from the category of \(k\)-Lie algebras to the
category of infinitesimal formal \(k\)-groups.}\footnote{Translator's note: the Gille--Polo re-edition (2011),
Expos\'e~VII\textsubscript{B}, observes here that, by 2.2.1, 1.3.6 and 2.7, an infinitesimal formal \(k\)-group
amounts to a cocommutative bialgebra over \(k\) that is connected (its natural filtration, 1.3.6, exhausts
it); so the corollary says that, over a field of characteristic~\(0\), \(\mathbf L\mapsto U(\mathbf L)\) and
\(H\mapsto\operatorname{Prim}(H)\) are quasi-inverse equivalences between Lie algebras and connected
cocommutative bialgebras. This was also obtained by B.~Kostant and, for bialgebras generated by their
primitive elements, by J.~W.~Milnor and J.~C.~Moore, \emph{On the structure of Hopf algebras}, Ann.\ of
Math.\ \textbf{81} (1965), 211--264. The re-edition adds that, for an Artinian ring \(k\) containing
\(\mathbb Q\), 3.2 and 3.3 give in the same way an equivalence \(\mathbf L\mapsto\mathscr G(\mathbf L)\)
from flat Lie algebras over \(\mathbf O_k\) to topologically flat infinitesimal formal \(k\)-groups.}

Indeed, as \(G\) ranges over the infinitesimal formal
\(k\)-groups, the functor
\[
  G\longmapsto\mathscr G\bigl(\operatorname{Lie}G\bigr)
\]
is isomorphic to the identity functor of the category of infinitesimal
groups \textup{(\SGARefLink{sga3:VIIB:3.3:cartier-theorem}{Theorem~3.3})}. Likewise, the functor
\[
  \mathbf L
  \longmapsto
  \operatorname{Lie}\bigl(\mathscr G(\mathbf L)\bigr)
\]
is isomorphic to the identity functor of the category of Lie algebras
\textup{(\SGARefLink{sga3:VIIB:3.1}{3.1})}.\footnote{Translator's note: by
\SGARefLink{sga3:VIIB:2.6:item:c}{2.6~c)}, \(\operatorname{Lie}(\mathscr G(\mathbf L))\) is the set of primitive elements
of the bialgebra \(U(\mathbf L)\) of \(\mathscr G(\mathbf L)\), and 3.1 identifies this set with \(\mathbf L\).
The Gille--Polo re-edition (2011), Expos\'e~VII\textsubscript{B}, makes this explicit by writing
\(\operatorname{Lie}(\mathscr G(\mathbf L))=\operatorname{Prim}U(\mathbf L)\).}
